\section*{Introducción}

La caracterización de los sensores RTD es fundamental para garantizar mediciones de temperatura precisas, repetibles y trazables en aplicaciones industriales y de instrumentación. Este proceso permite evaluar parámetros clave como la relación resistencia–temperatura, la linealidad, la sensibilidad y la estabilidad frente a variaciones ambientales y eléctricas, lo cual es crítico en sistemas de control de procesos, monitoreo térmico y adquisición de datos.

En la caracterización de sensores de temperatura, los PT100 y los NTC destacan por sus propiedades y aplicaciones complementarias: los PT100, como RTD normalizados, ofrecen alta precisión, excelente linealidad y estabilidad a largo plazo, lo que los hace ideales para instrumentación industrial, calibración y control de procesos donde se requiere trazabilidad y exactitud; por otro lado, los NTC presentan alta sensibilidad y rápida respuesta térmica, siendo más económicos y adecuados para aplicaciones de monitoreo, protección térmica y sistemas embebidos donde las variaciones de temperatura deben detectarse con rapidez, aunque con menor linealidad. La caracterización puntual de ambos sensores permite seleccionar el dispositivo adecuado y diseñar correctamente el acondicionamiento de señal según los requerimientos de precisión, rango y dinámica del sistema.